\DocumentMetadata{
  pdfversion=2.0,
  pdfstandard=ua-2,
  lang=en-US,
  tagging=on,
  tagging-setup={math/setup=mathml-SE,tabsorder=structure}
}
\documentclass[11pt,letterpaper]{article}
\usepackage[margin=0.68in,includefoot]{geometry}
\usepackage{newcomputermodern}
\usepackage{amsmath,amscd,mathtools}
\usepackage{tikz,adjustbox}
\providecommand{\square}{\mdlgwhtsquare}
\usepackage[hidelinks]{hyperref}
\hypersetup{
  pdftitle={Algebra First-Year Exam, May 2019, Part 2},
  pdfauthor={University of Florida Department of Mathematics},
  pdfsubject={Graduate mathematics examination},
  pdfdisplaydoctitle=true,
  bookmarksopen=true
}
\setcounter{secnumdepth}{0}
\setlength{\parindent}{0pt}
\setlength{\parskip}{0.28em}
\setlength{\emergencystretch}{3em}
\setlength{\leftmargini}{2.15em}
\setlength{\leftmarginii}{2.65em}
\setlength{\leftmarginiii}{3em}
\pagestyle{empty}
\raggedbottom
\allowdisplaybreaks
\NewTaggingSocketPlug{block/list/label}{examproblem}{%
  \tagstructbegin{tag=Lbl}\tagstructend%
  \tagstructbegin{tag=\LBody}%
  \tagstructbegin{tag=\ProblemHeadingTag,title-o={\ProblemHeadingText}}%
  \tagmcbegin{tag=\ProblemHeadingTag,actualtext={\ProblemHeadingText}}%
  #2%
  \tagmcend%
  \tagstructend%
}
\newenvironment{examproblems}[2]{%
  \def\ProblemHeadingTag{#2}%
  \AssignTaggingSocketPlug{block/list/label}{examproblem}%
  \begin{enumerate}%
  \setcounter{enumi}{#1}%
  \renewcommand{\labelenumi}{\textbf{\theenumi.}}%
  \setlength{\topsep}{0.35em}%
  \setlength{\partopsep}{0pt}%
  \setlength{\itemsep}{0.48em}%
  \setlength{\parsep}{0.18em}%
}{\end{enumerate}}
\newenvironment{examsubparts}{%
  \AssignTaggingSocketPlug{block/list/label}{default}%
  \begin{enumerate}%
  \setlength{\topsep}{0.18em}%
  \setlength{\partopsep}{0pt}%
  \setlength{\itemsep}{0.16em}%
  \setlength{\parsep}{0.1em}%
}{\end{enumerate}}
\newenvironment{examinstructions}{%
  \par\begingroup\itshape
}{%
  \par\endgroup\vspace{0.18em}
}
\makeatletter
\renewcommand\subsection{\@startsection{subsection}{2}{0pt}%
  {-1.25ex plus -.25ex minus -.1ex}%
  {0.55ex plus .1ex}%
  {\normalfont\large\bfseries}}
\renewcommand{\@maketitle}{%
  \newpage\null\vskip -1em%
  {\footnotesize\bfseries
    \href{https://www.ufl.edu/}{University of Florida}, \href{https://math.ufl.edu/}{Department of Mathematics}\par}%
  \vskip -0.5em%
  \tagstructbegin{tag=H1,title={Algebra First-Year Exam, May 2019, Part 2}}%
  \tagpdfparaOff%
  \tagmcbegin{tag=H1}%
    {\LARGE\bfseries\href{https://gma.math.ufl.edu/exams/algebra-first-year/}{Algebra First-Year Exam}, May 2019, Part 2\par}%
  \tagmcend%
  \tagpdfparaOn%
  \tagstructend%
  \vskip 0.25em%
  \tagmcbegin{artifact}%
    \hrule height 1pt%
  \tagmcend%
  \vskip 1em%
}
\makeatother
\title{Algebra First-Year Exam, May 2019, Part 2}
\author{}
\date{}
\begin{document}
\maketitle
\thispagestyle{empty}
\begin{examinstructions}
Answer four problems. (If you turn in more, the first four will be graded.) Within reason, you may quote theorems as long as you state them clearly.
\end{examinstructions}
\begin{examproblems}{0}{H2}
\def\ProblemHeadingText{Problem 1.}
\item
State and prove the Chinese Remainder Theorem for rings.
\def\ProblemHeadingText{Problem 2.}
\item
Let \(R=\mathbb{Z}[\sqrt{-1}]\) be the ring of Gaussian integers. Suppose that~\(F\) is a finite field of characteristic~\(5\), and \(\phi:R\to F\) is a (unital) ring homomorphism. Prove that \(R/\ker(\phi)\) has exactly~\(5\) elements.
\def\ProblemHeadingText{Problem 3.}
\item
Let \(\omega=\frac{-1+\sqrt{-3}}{2}\in\mathbb{C}\).
\begin{examsubparts}
\item[\textnormal{(a)}]
Prove that~\(\omega\) is a root of the polynomial \(f(x)=x^2+x+1\).
\item[\textnormal{(b)}]
Prove that the set \(\mathbb{Z}[\omega]=\{a+b\omega:a,b\in\mathbb{Z}\}\) is a subring of~\(\mathbb{C}\).
\item[\textnormal{(c)}]
Prove that \(\mathbb{Z}[\omega]\) is a Euclidean domain with respect to the norm \(N(\alpha)=\alpha\overline{\alpha}\) for all \(\alpha\in\mathbb{Z}[\omega]\), where \(\overline{\alpha}\) denotes the complex conjugate of~\(\alpha\).
\end{examsubparts}
\def\ProblemHeadingText{Problem 4.}
\item
Let~\(R\) be an integral domain, and~\(M\) be an~\(R\)-module.
\begin{examsubparts}
\item[\textnormal{(a)}]
Define the rank of~\(M\);
\item[\textnormal{(b)}]
Prove that if~\(M\) is a free~\(R\)-module of rank~\(n\), (where~\(n\) is a non-negative integer), then the rank of any submodule of~\(M\) is at most~\(n\).
\end{examsubparts}
\def\ProblemHeadingText{Problem 5.}
\item
This problem has two parts.
\begin{examsubparts}
\item[\textnormal{(a)}]
Let~\(p\) be an odd prime. Let \(A\in\operatorname{GL}(2,p)\) be an invertible two by two matrix over the field with~\(p\) elements such that~\(A\) has order~\(2\) and \(A\ne-I\) (where~\(I\) is the identity matrix). Prove that~\(A\) is conjugate to the diagonal matrix whose diagonal entries are \(-1\) and~\(1\).
\item[\textnormal{(b)}]
Give an example of an element of order~\(2\) in \(\operatorname{GL}(2,2)\).
\end{examsubparts}
\def\ProblemHeadingText{Problem 6.}
\item
Let \(F/K\) be a field extension.
\begin{examsubparts}
\item[\textnormal{(a)}]
Define what it means to say that \(\theta\in F\) is algebraic over~\(K\).
\item[\textnormal{(b)}]
Let~\(A\) be the set of all elements of~\(F\) that are algebraic over~\(K\). Prove that~\(A\) is a subfield of~\(F\).
\item[\textnormal{(c)}]
Prove that every element of~\(F\) that is not in~\(A\) is transcendental over~\(A\).
\end{examsubparts}
\end{examproblems}
\end{document}
